% Generated draft: requires mathematical and source review.
\documentclass[11pt]{article}
\usepackage[margin=25mm]{geometry}
\usepackage{fontspec}
\setmainfont{Latin Modern Roman}
\usepackage{amsmath,amssymb,mathtools,unicode-math}
\setmathfont{Latin Modern Math}
\usepackage{enumitem,needspace,xurl,hyperref}
\setlength{\parindent}{0pt}
\setlength{\parskip}{5pt}
\setlength{\emergencystretch}{3em}
\newcommand{\N}{\mathbb{N}}
\newcommand{\Z}{\mathbb{Z}}
\newcommand{\Q}{\mathbb{Q}}
\newcommand{\R}{\mathbb{R}}
\newcommand{\C}{\mathbb{C}}
\newcommand{\F}{\mathbb{F}}
\newcommand{\A}{\mathbb{A}}
\newcommand{\PP}{\mathbb{P}}
\newcommand{\E}{\mathbb{E}}
\newcommand{\eps}{\varepsilon}
\newcommand{\dd}{\,\mathrm{d}}
\begin{document}
\section*{UnsolvedMath ID: 30000229}
\textbf{Source problem code:} catalogue entry
\subsection{Definitions and mathematical statement}
Consider the two-dimensional Poisson problem in the report's energy space $V$ on a domain $\Omega\subset\R^2$, with right-hand side $f\in V^*$ and weak solution $u$ satisfying
\[
 \int_\Omega\nabla u\cdot\nabla v\,\dd x=\langle f,v\rangle
 \quad(v\in V).
\]
For example, in the usual homogeneous Dirichlet setting $V=H_0^1(\Omega)$ with energy norm $\|v\|_V=\|\nabla v\|_2$, and the dual norm is $H^{-1}$. The source's boundary and mesh conventions must be used for the actual problem.

Let $\mathbb T_n$ be its admissible partitions with at most $n$ elements, let $S^1(\mathcal T)\cap V$ be the conforming piecewise-linear space, and let $P^0(\mathcal T)$ be the piecewise-constant functions, viewed as functionals on $V$. Define
\[
 \sigma_n(u)=\inf_{\mathcal T\in\mathbb T_n}
              \inf_{v\in S^1(\mathcal T)\cap V}\|u-v\|_V,
 \qquad
 \delta_n(f)=\inf_{\mathcal T\in\mathbb T_n}
              \inf_{g\in P^0(\mathcal T)}\|f-g\|_{V^*}.
\]
For $0<s\le1/2$, membership in $\mathcal A^s$ means $\sigma_n(u)=O(n^{-s})$. The conjecture is
\[
 \forall s\in(0,1/2],\qquad
 u\in\mathcal A^s\quad\Longrightarrow\quad\delta_n(f)=O(n^{-s}).
\]
The constants may depend on the data and $s$ but not on $n$. The approximating partitions for $u$ and $f$ need not be the same. This is an existence implication; the original passage explicitly distinguishes existence of these data approximations from knowing how to construct them.

The extract does not specify domain/boundary regularity, admissible refinement and shape-regularity rules, the precise mesh-complexity offset, or all Sobolev hypotheses alluded to by its literature assessment. The notation above exposes these inputs without replacing the source mesh class by unrestricted partitions.

The literature assessment checked on 21 August 2026 reports optimal adaptive finite-element results conditional on the required data approximation property and positive results for important data classes, but no resolution of the broader implication over its full stated ranges. A theorem assuming that property does not prove that it follows from $u\in\mathcal A^s$. This is dated supplied partial-progress triage requiring independent verification.
\subsection{Short English statement}
If piecewise-linear adaptive meshes approximate a two-dimensional Poisson solution at rate $n^{-s}$ for $0<s\le1/2$, must piecewise constants approximate its forcing at the same rate in the dual energy norm? The question asks for existence, not an algorithm that finds those approximations.
\subsection{Sources}
\begin{itemize}
\item \emph{Schnelle L\"oser f\"ur partielle Differentialgleichungen} (2006), Oberwolfach Reports, DOI 10.4171/owr/2005/24, pp. 1299--1370, \url{https://doi.org/10.4171/owr/2005/24}; proposed year supplied as 2005.
\item Literature assessment checked 21 August 2026: original report, \url{https://doi.org/10.4171/owr/2005/24}; Stevenson, ``optimality of adaptive finite elements,'' \url{https://doi.org/10.1007/s10208-005-0183-0}.
\item Local UnsolvedMath registry, permanent ID $30000229$, exact source code \texttt{catalogue entry}, created 25 August 2026 and updated 16 September 2026. Supplied verification restores the Poisson/linear-element setting, $0<s\le1/2$, the energy rate and piecewise-constant dual-norm rate.
\end{itemize}
\end{document}
